% !TeX root = ../main.tex
\newcommand{\PreprocessAgeMedian}{52}
\newcommand{\PreprocessAgeChanged}{1}
\newcommand{\PreprocessTrainN}{90,000}
\newcommand{\PreprocessValidationN}{30,000}
\newcommand{\PreprocessTestN}{30,000}


按家属人数、逾期编码、收入状态依次处理，标记允许重叠，规则见表~\ref{tab:cleaning-rules}。
相较于无标记的直接填补，保留指示或独立状态可避免混同图~\ref{fig:training-feature-risk} 中风险不同的人群；
本节选定数据准备方案，预测性能仍由后续验证比较。

\begin{table}[!htbp]
  \centering
  \caption{缺失与特殊值的冻结处理规则}
  \label{tab:cleaning-rules}
  \footnotesize
  \setlength{\tabcolsep}{3pt}
  \renewcommand{\arraystretch}{1.04}
  \begin{tabular}{p{2.0cm}p{13.1cm}}
    \toprule
    对象 & 处理及迁移规则 \\
    \midrule
    家属人数 & 原缺失先标记为 1 再补 0，原生零值标记为 0，保留两类差异。 \\
    三类逾期 & 保留原值，逐列标记 96/98 并合为该变量的特殊档；共同标记仅追踪任一列命中，避免重复计人数。 \\
    月收入 & 保留原值；缺失、$[0,10)$、$[10,+\infty)$ 分别为 \texttt{MISSING}、\texttt{LOW}、\texttt{REGULAR}。
    前两者路由特殊档，10 归普通档，原 NA 由缺失状态映射。 \\
    年龄 & 缺失或不大于 0 时补训练集有效年龄中位数 \PreprocessAgeMedian{} 岁并标记。
    本次改动 \PreprocessAgeChanged{} 条；训练中恒定的异常标记仅供审计，不估计系数。 \\
    比例与长尾 & DebtRatio 在各收入状态下保留；比例大于 1、高收入及其他合法长尾均不截尾或删行，普通分箱和 WOE 留待模块三学习。 \\
    \bottomrule
  \end{tabular}
\end{table}

\begin{samepage}
\paragraph{家属人数缺失的处理依据}
在完整训练集中，家属人数缺失组的坏样本率为 \DependentsMissingRisk，原生零值组为
\DependentsZeroRisk，均低于家属人数为正组的 8.08\%，呈现相对较低的边际风险。
因此，选择 0 作为缺失记录的数值占位，同时保留原缺失指示，使其在数值表示上与零值组衔接，
并允许后续模型识别缺失状态的附加信息。这一处理不将两组风险视为相同，
也不认定缺失客户的真实家属人数为 0；相较于直接补 0，缺失标记保留了两类记录的区别。\par
\end{samepage}

\paragraph{逾期特殊编码的保留依据}
三类逾期字段的 96/98 在训练集中命中同一批 \LateSpecialN{} 条记录，其中 91 条为坏样本，
坏样本率达 \LateSpecialRisk，明显高于训练集总体的 6.68\%。
该编码状态具有较强的样本风险区分信息，直接删除会丢失目标事件高度集中的记录。
鉴于公开字典未明确其业务含义，本研究保留原始记录，将 96/98 合并为特殊档，
既保留其风险信息，也避免将未知编码当作正常逾期次数排序。
其中 96 仅有 2 条训练记录，与 98 合并可避免对极稀疏子组单独估计风险；
三列共享的特殊状态在联合建模时仍需检查重复信息。
\paragraph{收入三种状态的划分依据}
为考察非逾期特殊编码人群中的收入关系，在不含 96/98 的 89,839 条训练记录中比较：
收入缺失、$[0,10)$ 及 $[10,+\infty)$ 三组分别有 17,768、1,347、70,724 条，
坏样本率分别为 5.57\%、3.34\%、6.91\%。
缺失与极低收入组不仅区别于普通收入组，彼此也呈现不同的边际风险，
因此分别设置 \texttt{MISSING} 和 \texttt{LOW} 状态，避免预先合并，
普通收入组再于后续学习常规分箱。这是基于记录状态与风险分布的编码选择，
并不意味着收入越低越安全。上述条件样本仅用于解释收入分层，实际转换仍覆盖所有记录，
包括命中 96/98 的客户，收入状态与逾期标记可以重叠。

数据处理后，三个子集的样本个数不变。
参数只从训练特征估计，冻结后统一用于验证、测试及新客户；交叉验证在折内重估。
合法超训练范围值保留，由后续开口箱覆盖；缺字段、非法格式、无穷值、负收入或比例、
负或非整数计数及未约定缺失均报字段级错误，不静默删行。
新客户无需标签；完整规则清单、变动人数和标记重叠记录随代码保存。
